% --- Archon physics formalization source begin ---
% archon:physics
% archon:covers PhyXMiniProblems/problem_phyx_mini_0505.lean
% archon:source-report reports/phyx_mini/problem_phyx_mini_0505.source.json

\chapter{Physics problem phyx\_mini\_0505}
\label{ch:PhyXMiniProblems_problem_phyx_mini_0505}

\paragraph{Problem source.}
\#\# Physical scenario

Figure shows \textbackslash{}( |\textbackslash{}psi(x)|\textasciicircum{}2 \textbackslash{}) for the electrons in an experiment.

\#\# Question

If \textbackslash{}( 10\textasciicircum{}4 \textbackslash{}) electrons are detected, how many are expected to land in the interval \textbackslash{}( -0.30 \textbackslash{}text\{ cm\} \textbackslash{}leq x \textbackslash{}leq 0.30 \textbackslash{}text\{ cm\} \textbackslash{})?

\#\# Answer choices

- A: A: 500

- B: B: 1000

- C: C: 900

- D: D: 1300

\#\# Figure caption (auxiliary; use the image as primary evidence)

The image is a graph plotting \textbackslash{}(|\textbackslash{}psi(x)|\textasciicircum{}2\textbackslash{}) against \textbackslash{}(x\textbackslash{}).

- **Axes**: 
  - The x-axis represents \textbackslash{}(x\textbackslash{}) in centimeters (cm) and ranges from -1 to 1.
  - The y-axis represents \textbackslash{}(|\textbackslash{}psi(x)|\textasciicircum{}2\textbackslash{}) in inverse centimeters (cm\textbackslash{}(\textasciicircum{}\{-1\}\textbackslash{})) and ranges from 0 to 1.

- **Plot**: 
  - The function forms a triangular shape.
  - At \textbackslash{}(x = 0\textbackslash{}), \textbackslash{}(|\textbackslash{}psi(x)|\textasciicircum{}2 = 0\textbackslash{}).
  - At \textbackslash{}(x = -1\textbackslash{}) and \textbackslash{}(x = 1\textbackslash{}), \textbackslash{}(|\textbackslash{}psi(x)|\textasciicircum{}2 = 1\textbackslash{}).
  - The lines connecting these points are straight, indicating a linear relationship.

- **Text**:
  - The y-axis is labeled as \textbackslash{}(|\textbackslash{}psi(x)|\textasciicircum{}2\textbackslash{}) (cm\textbackslash{}(\textasciicircum{}\{-1\}\textbackslash{})).
  - The x-axis is labeled as \textbackslash{}(x\textbackslash{}) (cm).

\#\# Dataset metadata

Quantum Phenomena; Physical Model Grounding Reasoning, Numerical Reasoning

\paragraph{Recorded answer/context.}
C: C: 900

\paragraph{Figure/image path.}
/root/proposal\_for\_physic/science-mango/phyx\_mini\_run/phyx\_data/test\_image/505.png

\paragraph{Formalization target.}
create a compiling Lean file with sorry bodies at `PhyXMiniProblems/problem\_phyx\_mini\_0505.lean`. The Lean declarations must preserve the physical quantities, dimensions or dimensional roles, figure labels, governing-law hypotheses, and final relation expressed by this problem.
Use Mathlib/Physlib names found through LeanExplore where available. If a physics API is missing, introduce faithful local abstractions rather than scalar placeholder aliases.

\begin{theorem}[Physics formalization target]
\label{thm:physics:phyx_mini_0505:target}
\lean{PhyXMiniProblems.ProblemPhyXMini0505.problem_phyx_mini_0505}
\uses{def:physics:phyx-mini-0505:phyxminiproblems-problemphyxmini0505-electronpositionexperiment, def:physics:phyx-mini-0505:phyxminiproblems-problemphyxmini0505-probabilitydensityfigureaxes, def:physics:phyx-mini-0505:phyxminiproblems-problemphyxmini0505-detectionquestionreadout, def:physics:phyx-mini-0505:phyxminiproblems-problemphyxmini0505-matchessuppliedprobabilitydensityfigure, def:physics:phyx-mini-0505:phyxminiproblems-problemphyxmini0505-matchesdetectionquestion, def:physics:phyx-mini-0505:phyxminiproblems-problemphyxmini0505-satisfiesbornandcountinglaws, def:physics:phyx-mini-0505:phyxminiproblems-problemphyxmini0505-displayedexpectedcount}
The assigned autoformalize agent should translate this physics problem into Lean declarations in the covered file, with theorem and lemma proof bodies written as `by sorry`.
\end{theorem}
\begin{proof}
This is an autoformalization task, not a proof task. Produce statements that can later be proved by the physics prover without weakening the physical model.
\end{proof}

% --- Archon declaration topology begin ---
\section{Declaration topology}
\begin{definition}[position Axis Unit]
  \label{def:physics:phyx-mini-0505:phyxminiproblems-problemphyxmini0505-positionaxisunit}
  \lean{PhyXMiniProblems.ProblemPhyXMini0505.positionAxisUnit}
  The physical length unit used for every position-axis scalar readout
\end{definition}
\begin{proof}
  This is the declared model definition of position Axis Unit.
\end{proof}
\begin{definition}[probability Density Dimension]
  \label{def:physics:phyx-mini-0505:phyxminiproblems-problemphyxmini0505-probabilitydensitydimension}
  \lean{PhyXMiniProblems.ProblemPhyXMini0505.probabilityDensityDimension}
  The physical dimension carried by the graph's vertical-axis readouts
\end{definition}
\begin{proof}
  This is the declared model definition of probability Density Dimension.
\end{proof}
\begin{definition}[Electron Position Experiment]
  \label{def:physics:phyx-mini-0505:phyxminiproblems-problemphyxmini0505-electronpositionexperiment}
  \lean{PhyXMiniProblems.ProblemPhyXMini0505.ElectronPositionExperiment}
  The physical experiment, represented through calibrated scalar readouts. The argument of waveFunctionInCentimeterCoordinates and probabilityDensityPerCentimeter is the numerical coordinate in centimeters. The complex wavefunction readout therefore has the implicit unit cm⁻¹ᐟ², and its squared modulus has unit cm⁻¹. Interval probabilities are dimensionless. An expected count is real-valued even though the number of detected electrons is natural-valued
\end{definition}
\begin{proof}
  This is the declared model definition of Electron Position Experiment.
\end{proof}
\begin{definition}[Probability Density Figure Axes]
  \label{def:physics:phyx-mini-0505:phyxminiproblems-problemphyxmini0505-probabilitydensityfigureaxes}
  \lean{PhyXMiniProblems.ProblemPhyXMini0505.ProbabilityDensityFigureAxes}
  Numerical labels on the position-probability-density plot
\end{definition}
\begin{proof}
  This is the declared model definition of Probability Density Figure Axes.
\end{proof}
\begin{definition}[Detection Question Readout]
  \label{def:physics:phyx-mini-0505:phyxminiproblems-problemphyxmini0505-detectionquestionreadout}
  \lean{PhyXMiniProblems.ProblemPhyXMini0505.DetectionQuestionReadout}
  Numerical centimeter coordinates defining the interval asked about
\end{definition}
\begin{proof}
  This is the declared model definition of Detection Question Readout.
\end{proof}
\begin{definition}[Matches Supplied Probability Density Figure]
  \label{def:physics:phyx-mini-0505:phyxminiproblems-problemphyxmini0505-matchessuppliedprobabilitydensityfigure}
  \lean{PhyXMiniProblems.ProblemPhyXMini0505.MatchesSuppliedProbabilityDensityFigure}
  \uses{def:physics:phyx-mini-0505:phyxminiproblems-problemphyxmini0505-electronpositionexperiment, def:physics:phyx-mini-0505:phyxminiproblems-problemphyxmini0505-probabilitydensityfigureaxes}
  The supplied figure has endpoints (-1 cm, 1 cm⁻¹) and (1 cm, 1 cm⁻¹), a node at (0 cm, 0 cm⁻¹), straight segments between those points, and zero density outside [-1 cm, 1 cm]
\end{definition}
\begin{proof}
  This is the declared model definition of Matches Supplied Probability Density Figure.
\end{proof}
\begin{definition}[Matches Detection Question]
  \label{def:physics:phyx-mini-0505:phyxminiproblems-problemphyxmini0505-matchesdetectionquestion}
  \lean{PhyXMiniProblems.ProblemPhyXMini0505.MatchesDetectionQuestion}
  \uses{def:physics:phyx-mini-0505:phyxminiproblems-problemphyxmini0505-electronpositionexperiment, def:physics:phyx-mini-0505:phyxminiproblems-problemphyxmini0505-detectionquestionreadout}
  The prose data: 10⁴ detections and the interval [-0.30, 0.30] cm
\end{definition}
\begin{proof}
  This is the declared model definition of Matches Detection Question.
\end{proof}
\begin{definition}[Satisfies Born And Counting Laws]
  \label{def:physics:phyx-mini-0505:phyxminiproblems-problemphyxmini0505-satisfiesbornandcountinglaws}
  \lean{PhyXMiniProblems.ProblemPhyXMini0505.SatisfiesBornAndCountingLaws}
  \uses{def:physics:phyx-mini-0505:phyxminiproblems-problemphyxmini0505-electronpositionexperiment}
  General physical laws used in the calculation. The density is the squared modulus of a square-integrable, normalized one-dimensional wavefunction. The Born rule obtains interval probability by integrating that density, and the expected landing count is total detections times interval probability
\end{definition}
\begin{proof}
  This is the declared model definition of Satisfies Born And Counting Laws.
\end{proof}
\begin{definition}[Answer Choice]
  \label{def:physics:phyx-mini-0505:phyxminiproblems-problemphyxmini0505-answerchoice}
  \lean{PhyXMiniProblems.ProblemPhyXMini0505.AnswerChoice}
  Labels of the four displayed multiple-choice answers
\end{definition}
\begin{proof}
  This is the declared model definition of Answer Choice.
\end{proof}
\begin{definition}[displayed Expected Count]
  \label{def:physics:phyx-mini-0505:phyxminiproblems-problemphyxmini0505-displayedexpectedcount}
  \lean{PhyXMiniProblems.ProblemPhyXMini0505.displayedExpectedCount}
  \uses{def:physics:phyx-mini-0505:phyxminiproblems-problemphyxmini0505-answerchoice}
  Expected-electron counts printed beside the displayed answer choices
\end{definition}
\begin{proof}
  This is the declared model definition of displayed Expected Count.
\end{proof}
\begin{definition}[recorded Dataset Answer]
  \label{def:physics:phyx-mini-0505:phyxminiproblems-problemphyxmini0505-recordeddatasetanswer}
  \lean{PhyXMiniProblems.ProblemPhyXMini0505.recordedDatasetAnswer}
  \uses{def:physics:phyx-mini-0505:phyxminiproblems-problemphyxmini0505-answerchoice}
  The answer label recorded in the source dataset, retained as metadata
\end{definition}
\begin{proof}
  This is the declared model definition of recorded Dataset Answer.
\end{proof}
\begin{lemma}[queried Interval Probability eq nine Hundredths]
  \label{lem:physics:phyx-mini-0505:phyxminiproblems-problemphyxmini0505-queriedintervalprobability-eq-ninehundredths}
  \lean{PhyXMiniProblems.ProblemPhyXMini0505.queriedIntervalProbability_eq_nineHundredths}
  \uses{def:physics:phyx-mini-0505:phyxminiproblems-problemphyxmini0505-electronpositionexperiment, def:physics:phyx-mini-0505:phyxminiproblems-problemphyxmini0505-probabilitydensityfigureaxes, def:physics:phyx-mini-0505:phyxminiproblems-problemphyxmini0505-detectionquestionreadout, def:physics:phyx-mini-0505:phyxminiproblems-problemphyxmini0505-matchessuppliedprobabilitydensityfigure, def:physics:phyx-mini-0505:phyxminiproblems-problemphyxmini0505-matchesdetectionquestion, def:physics:phyx-mini-0505:phyxminiproblems-problemphyxmini0505-satisfiesbornandcountinglaws}
  The area under the plotted density over [-0.30, 0.30] cm is 0.09
\end{lemma}
\begin{proof}
  The conclusion follows from the governing relations and figure readouts named in the statement of queried Interval Probability eq nine Hundredths.
\end{proof}
% --- Archon declaration topology end ---
% --- Archon physics formalization source end ---
